% vkr.tex — Шаблон ВКР, Университет ИТМО
% ГОСТ 7.32-2017 / ЛНАОБУЧ-СМК-03-05-2022
\documentclass[a4paper,14pt]{extarticle}

% preamble.tex — Преамбула шаблона ВКР ИТМО
% Оформление по ГОСТ 7.32-2017 и ЛНАОБУЧ-СМК-03-05-2022

\usepackage[T2A]{fontenc}
\usepackage[utf8]{inputenc}
\usepackage[english,russian]{babel}

\usepackage{mathptmx}                        % Times New Roman

\usepackage[nodisplayskipstretch]{setspace}   % Полуторный интервал
\onehalfspacing

% Поля: лев 30, прав 15, верх 20, низ 20 мм (п. 4.2)
\usepackage[top=20mm, bottom=20mm, left=30mm, right=15mm]{geometry}

\parindent=1.25cm
\usepackage{indentfirst}

% Нумерация страниц по центру внизу
\usepackage{fancyhdr}
\pagestyle{fancy}
\fancyhf{}
\fancyfoot[C]{\thepage}
\renewcommand{\headrulewidth}{0pt}
\renewcommand{\footrulewidth}{0pt}

% Содержание
\usepackage{titletoc}
\usepackage{tocloft}
\usepackage{etoolbox}
\setcounter{tocdepth}{3}

\makeatletter
\patchcmd{\@cftmaketoctitle}{\cfttoctitlefont\contentsname}{\cfttoctitlefont\MakeUppercase{\contentsname}}{}{}
\makeatother

\renewcommand{\cftsecleader}{\cftdotfill{\cftdotsep}}

% Заголовки (п. 4.4.4): с абзацного отступа, полужирный
\usepackage{titlesec}
\usepackage{textcase}

\titleformat{\section}[block]{\hspace{\parindent}\bfseries\normalsize}
                         {\arabic{section}}{1ex}{\MakeUppercase}
\titlespacing*{\section}{0pt}{42pt}{42pt}

\titleformat{\subsection}[block]{\hspace{\parindent}\bfseries\normalsize}
                         {\arabic{section}.\arabic{subsection}}{1ex}{}
\titlespacing*{\subsection}{0pt}{\baselineskip}{\baselineskip}

\titleformat{\subsubsection}[block]{\hspace{\parindent}\bfseries\normalsize}
                         {\arabic{section}.\arabic{subsection}.\arabic{subsubsection}}{1ex}{}
\titlespacing*{\subsubsection}{0pt}{\baselineskip}{\baselineskip}

% Подписи: разделитель — короткое тире (п. 4.5.3, 4.6.3)
\usepackage{caption}
\DeclareCaptionLabelSeparator{endash}{ -- }
\captionsetup{labelsep=endash}
\captionsetup[figure]{justification=centering, font=normalsize}
\captionsetup[table]{singlelinecheck=false, justification=raggedright, font=normalsize}

\addto\captionsrussian{
  \def\figurename{Рисунок}
  \def\tablename{Таблица}
  \def\contentsname{Содержание}
  \def\refname{СПИСОК ИСПОЛЬЗОВАННЫХ ИСТОЧНИКОВ}
}

% Сквозная нумерация
\renewcommand{\theequation}{\arabic{equation}}
\renewcommand{\thetable}{\arabic{table}}
\renewcommand{\thefigure}{\arabic{figure}}

\usepackage{amsmath}
\usepackage{amssymb}
\usepackage{amsthm}
\usepackage{amsfonts}

\usepackage[pdftex]{graphicx}
\graphicspath{{images/}}
\usepackage{float}

\usepackage{multirow}
\usepackage{longtable}
\usepackage{booktabs}
\usepackage{array}

% Списки: маркер — тире
\usepackage{enumitem}
\setlist[enumerate]{topsep=0pt, itemsep=0ex, partopsep=1ex, parsep=1ex}
\setlist[itemize]{topsep=0pt, itemsep=0ex, partopsep=1ex, parsep=1ex}
\def\labelitemi{--}

% Листинги кода с поддержкой кириллицы
\usepackage{listings}
\usepackage{xcolor}
\lstset{
  basicstyle=\small\ttfamily,
  breaklines=true,
  frame=single,
  xleftmargin=\parindent,
  showstringspaces=false,
  extendedchars=true,
  literate=
    {а}{{\cyra}}1 {б}{{\cyrb}}1 {в}{{\cyrv}}1 {г}{{\cyrg}}1
    {д}{{\cyrd}}1 {е}{{\cyre}}1 {ё}{{\cyryo}}1 {ж}{{\cyrzh}}1
    {з}{{\cyrz}}1 {и}{{\cyri}}1 {й}{{\cyrishrt}}1 {к}{{\cyrk}}1
    {л}{{\cyrl}}1 {м}{{\cyrm}}1 {н}{{\cyrn}}1 {о}{{\cyro}}1
    {п}{{\cyrp}}1 {р}{{\cyrr}}1 {с}{{\cyrs}}1 {т}{{\cyrt}}1
    {у}{{\cyru}}1 {ф}{{\cyrf}}1 {х}{{\cyrh}}1 {ц}{{\cyrc}}1
    {ч}{{\cyrch}}1 {ш}{{\cyrsh}}1 {щ}{{\cyrshch}}1 {ъ}{{\cyrhrdsn}}1
    {ы}{{\cyrery}}1 {ь}{{\cyrsftsn}}1 {э}{{\cyrerev}}1 {ю}{{\cyryu}}1
    {я}{{\cyrya}}1
    {А}{{\CYRA}}1 {Б}{{\CYRB}}1 {В}{{\CYRV}}1 {Г}{{\CYRG}}1
    {Д}{{\CYRD}}1 {Е}{{\CYRE}}1 {Ё}{{\CYRYO}}1 {Ж}{{\CYRZH}}1
    {З}{{\CYRZ}}1 {И}{{\CYRI}}1 {Й}{{\CYRISHRT}}1 {К}{{\CYRK}}1
    {Л}{{\CYRL}}1 {М}{{\CYRM}}1 {Н}{{\CYRN}}1 {О}{{\CYRO}}1
    {П}{{\CYRP}}1 {Р}{{\CYRR}}1 {С}{{\CYRS}}1 {Т}{{\CYRT}}1
    {У}{{\CYRU}}1 {Ф}{{\CYRF}}1 {Х}{{\CYRH}}1 {Ц}{{\CYRC}}1
    {Ч}{{\CYRCH}}1 {Ш}{{\CYRSH}}1 {Щ}{{\CYRSHCH}}1 {Ъ}{{\CYRHRDSN}}1
    {Ы}{{\CYRERY}}1 {Ь}{{\CYRSFTSN}}1 {Э}{{\CYREREV}}1 {Ю}{{\CYRYU}}1
    {Я}{{\CYRYA}}1,
}

\usepackage[pdftex]{hyperref}
\hypersetup{unicode=true, colorlinks=false, pdfborder={0 0 0}}
\expandafter\def\expandafter\UrlBreaks\expandafter{\UrlBreaks
  \do\a\do\b\do\c\do\d\do\e\do\f\do\g\do\h\do\i\do\j
  \do\k\do\l\do\m\do\n\do\o\do\p\do\q\do\r\do\s\do\t
  \do\u\do\v\do\w\do\x\do\y\do\z\do\A\do\B\do\C\do\D
  \do\E\do\F\do\G\do\H\do\I\do\J\do\K\do\L\do\M\do\N
  \do\O\do\P\do\Q\do\R\do\S\do\T\do\U\do\V\do\W\do\X
  \do\Y\do\Z}

\usepackage{cmap}

\usepackage[numbers,sort&compress]{natbib}
\makeatletter
\renewcommand\@biblabel[1]{#1.}
\makeatother

\usepackage{pdflscape}

\tolerance 1414
\hbadness 1414
\emergencystretch 1.5em
\hfuzz 0.3pt
\widowpenalty=10000
\vfuzz \hfuzz
\raggedbottom


\begin{document}

% Титульный лист, задание и аннотация оформляются в ИСУ (п. 3.3)
% \begin{titlepage}
% \thispagestyle{empty}
% \begin{center}
% \textbf{МИНИСТЕРСТВО НАУКИ И ВЫСШЕГО ОБРАЗОВАНИЯ РОССИЙСКОЙ ФЕДЕРАЦИИ}
% \vspace{0.5cm}
% \textbf{ФЕДЕРАЛЬНОЕ ГОСУДАРСТВЕННОЕ АВТОНОМНОЕ ОБРАЗОВАТЕЛЬНОЕ\\
% УЧРЕЖДЕНИЕ ВЫСШЕГО ОБРАЗОВАНИЯ\\
% <<НАЦИОНАЛЬНЫЙ ИССЛЕДОВАТЕЛЬСКИЙ УНИВЕРСИТЕТ ИТМО>>}
% \vspace{1cm}
% Факультет \underline{\hspace{6cm}}
% \vspace{0.3cm}
% Направление подготовки \underline{\hspace{6cm}}
% \vspace{0.3cm}
% Квалификация \underline{\hspace{6cm}}
% \vspace{2cm}
% {\large \textbf{ВЫПУСКНАЯ КВАЛИФИКАЦИОННАЯ РАБОТА}}
% \vspace{0.5cm}
% на тему:
% \vspace{0.3cm}
% <<\underline{\hspace{10cm}}>>
% \end{center}
% \vspace{2cm}
% \begin{flushright}
% \begin{minipage}{0.5\textwidth}
% Выполнил(а) студент(ка) группы \underline{\hspace{2cm}}\\[0.5cm]
% \underline{\hspace{\textwidth}}\\
% {\small (Фамилия И.О.)}\\[0.5cm]
% Руководитель ВКР\\[0.3cm]
% \underline{\hspace{\textwidth}}\\
% {\small (должность, учёная степень, звание)}\\[0.3cm]
% \underline{\hspace{\textwidth}}\\
% {\small (Фамилия И.О.)}\\[0.5cm]
% <<Допускается к защите>>\\[0.3cm]
% Руководитель ОП\\[0.3cm]
% \underline{\hspace{\textwidth}}\\
% {\small (Фамилия И.О., подпись)}
% \end{minipage}
% \end{flushright}
% \vfill
% \begin{center}
% Санкт-Петербург\\
% \the\year
% \end{center}
% \end{titlepage}

% === СОДЕРЖАНИЕ ===
\setcounter{page}{4}
\thispagestyle{fancy}
\tableofcontents
\newpage

% === СПИСОК СОКРАЩЕНИЙ И УСЛОВНЫХ ОБОЗНАЧЕНИЙ (п. 4.8) ===
\phantomsection
\begin{center}
\textbf{СПИСОК СОКРАЩЕНИЙ И УСЛОВНЫХ ОБОЗНАЧЕНИЙ}
\end{center}
\addcontentsline{toc}{section}{СПИСОК СОКРАЩЕНИЙ И УСЛОВНЫХ ОБОЗНАЧЕНИЙ}

\noindent
ВКР -- выпускная квалификационная работа\\
ГИА -- государственная итоговая аттестация\\
ГЭК -- государственная экзаменационная комиссия\\
ПО -- программное обеспечение

\newpage

% === ТЕРМИНЫ И ОПРЕДЕЛЕНИЯ (п. 4.9) ===
\phantomsection
\begin{center}
\textbf{ТЕРМИНЫ И ОПРЕДЕЛЕНИЯ}
\end{center}
\addcontentsline{toc}{section}{ТЕРМИНЫ И ОПРЕДЕЛЕНИЯ}

\noindent
Алгоритм -- конечная последовательность точно определённых действий для решения задачи\\
База данных -- организованная совокупность данных, предназначенная для хранения, обработки и извлечения информации\\
Интерфейс -- совокупность средств и методов взаимодействия между элементами системы

\newpage

% === ВВЕДЕНИЕ ===
\phantomsection
\begin{center}
\textbf{ВВЕДЕНИЕ}
\end{center}
\addcontentsline{toc}{section}{ВВЕДЕНИЕ}

Во введении.

\newpage

% === ГЛАВА 1 ===
\section{Название первой главы}

Текст главы

\subsection{Подраздел первой главы}

Текст подраздела.

\subsection{Ещё один подраздел}

Текст ещё одного подраздела.

\subsubsection{Пункт подраздела}

При необходимости подразделы могут содержать пункты с нумерацией вида 1.2.1, 1.2.2 и~т.\,д.

\newpage

% === ГЛАВА 2 ===
\section{Название второй главы}

Текст второй главы

\subsection{Пример ссылки на рисунок}

Иллюстрации нумеруются сквозной нумерацией (п.~4.5.3). На все иллюстрации должны быть ссылки в тексте (п.~4.5.4).

Схема архитектуры системы представлена на рисунке~\ref{fig:example}.

\begin{figure}[H]
  \centering
  % \includegraphics[width=0.8\textwidth]{имя_файла}
  \rule{0.6\textwidth}{4cm} % заглушка
  \caption{Пример рисунка}
  \label{fig:example}
\end{figure}

\subsection{Пример таблицы}

Таблицы нумеруются сквозной нумерацией (п.~4.6.2). На все таблицы должны быть ссылки в тексте.

Результаты эксперимента представлены в таблице~\ref{tab:results}.

\begin{table}[H]
\caption{Результаты экспериментального исследования}
\label{tab:results}
\centering
\begin{tabular}{|c|c|c|c|}
\hline
\textbf{Параметр} & \textbf{Значение 1} & \textbf{Значение 2} & \textbf{Значение 3} \\ \hline
Метрика A          & 0.95                & 0.87                & 0.91                \\ \hline
Метрика B          & 0.78                & 0.82                & 0.80                \\ \hline
Метрика C          & 0.64                & 0.71                & 0.68                \\ \hline
\end{tabular}
\end{table}

\newpage

% === ГЛАВА 3 ===
\section{Название третьей главы}

Текст третьей главы

\subsection{Пример маркированного списка}

\begin{itemize}
  \item первый пункт;
  \item второй пункт;
  \item третий пункт.
\end{itemize}

\subsection{Пример листинга кода}

Пример программного кода

\begin{lstlisting}[language=Python, caption={Пример функции на Python}]
def hello(name):
    """Приветствие пользователя."""
    print(f"Привет, {name}!")

if __name__ == "__main__":
    hello("мир")
\end{lstlisting}

\newpage

% === ЗАКЛЮЧЕНИЕ ===
\phantomsection
\begin{center}
\textbf{ЗАКЛЮЧЕНИЕ}
\end{center}
\addcontentsline{toc}{section}{ЗАКЛЮЧЕНИЕ}

В заключении ВКР обобщаются результаты выполненного исследования и делаются выводы, даются рекомендации, определяются перспективы дальнейшей разработки темы.

\newpage

% === СПИСОК ИСПОЛЬЗОВАННЫХ ИСТОЧНИКОВ ===
\phantomsection
\addcontentsline{toc}{section}{СПИСОК ИСПОЛЬЗОВАННЫХ ИСТОЧНИКОВ}
\bibliographystyle{unsrtnat}
\nocite{*}
\bibliography{references}

\newpage

% === ПРИЛОЖЕНИЕ А (п. 4.11) ===
\phantomsection
\begin{center}
ПРИЛОЖЕНИЕ А
\end{center}
\addcontentsline{toc}{section}{ПРИЛОЖЕНИЕ А}
\label{app:a}

\begin{center}
\textbf{Название приложения А}
\end{center}

Содержимое приложения А.

\newpage

% === ПРИЛОЖЕНИЕ Б ===
\phantomsection
\begin{center}
ПРИЛОЖЕНИЕ Б
\end{center}
\addcontentsline{toc}{section}{ПРИЛОЖЕНИЕ Б}
\label{app:b}

\begin{center}
\textbf{Название приложения Б}
\end{center}

Содержимое приложения Б.

\end{document}
